\chapter{Die Legende der Simulation}

\section{Einleitung}

Die Simulation der IWT erzeugt eine strukturierte Ausgabe, die den Zustand des Systems zu jedem Zeitschritt beschreibt.

Diese Ausgabe ist nicht nur ein technisches Protokoll. Sie ist der Beweis, dass die Theorie funktioniert.

Dieses Kapitel erklärt die Ausgabe der Simulation und ihre physikalische Bedeutung.

\section{Die Kernstatistik}

Die erste Zeile der Ausgabe zeigt die Kernstatistik des Systems:

\begin{verbatim}
#: 499                       - Aktuelle Iteration (Zeitschritt)
|Q_max|: 240806.394          - Maximalwert des Bohm-Potentials Q
I_total: 2.955               - Summe der Amplituden |I| über alle Knoten
I_min: 0.000                 - Minimale Amplitude |I| (Vakuum)
I_max: 0.063                 - Maximale Amplitude |I| (Struktur)
Deviation: 6.214             - Abweichung von der Anfangsbedingung
\end{verbatim}

Die Bedeutung:

\begin{itemize}
\item \textbf{Iteration:} Der aktuelle Zeitschritt. Die Simulation läuft stabil.
\item \textbf{\( |Q_{\text{max}}| \):} Die maximale Stärke des Bohm-Potentials. Ein großer Wert zeigt an, dass \( Q \) aktiv ist und Strukturen organisiert.
\item \textbf{\( I_{\text{total}} \):} Die Summe der Amplituden. Ein Maß für die gesamte Informationsmenge.
\item \textbf{\( I_{\text{min}} \):} Die minimale Amplitude. Das Vakuum.
\item \textbf{\( I_{\text{max}} \):} Die maximale Amplitude. Die Struktur.
\item \textbf{Deviation:} Die Abweichung von der Anfangsbedingung. Sie ist klein – das System ist stabil.
\end{itemize}

\section{Die Energiestatistik}

Die zweite Zeile zeigt die Energiestatistik:

\begin{verbatim}
Summe |I|^2: 2.955           - Summe der Betragsquadrate |I|^2
Info_Dev: 6.214              - Relative Abweichung von Summe |I|^2
\end{verbatim}

Die Bedeutung:

\begin{itemize}
\item \textbf{\( \sum |I|^2 \):} Die Summe der Betragsquadrate. Sie ist konstant – Axiom 2 ist erfüllt.
\item \textbf{Info\_Dev:} Die relative Abweichung von \( \sum |I|^2 \). Sie ist nahe Null – die Information bleibt erhalten.
\end{itemize}

\section{Die Fluktuationsstatistik}

Die dritte Zeile zeigt die Fluktuationsstatistik:

\begin{verbatim}
xi_real_mean: -0.003         - Mittelwert der Realteil-Fluktuationen
xi_imag_mean: -0.003         - Mittelwert der Imaginärteil-Fluktuationen
xi_real_sigma: 0.004         - Standardabweichung der Realteil-Fluktuationen
xi_imag_sigma: 0.004         - Standardabweichung der Imaginärteil-Fluktuationen
scale: 0.707                 - Skalierungsfaktor der Fluktuationen
\end{verbatim}

Die Bedeutung:

\begin{itemize}
\item \textbf{Mittelwerte:} Die Fluktuationen sind symmetrisch um Null.
\item \textbf{Standardabweichungen:} Die Fluktuationen sind isotrop – es gibt keine Vorzugsrichtung.
\item \textbf{Scale:} Der Skalierungsfaktor \( \sqrt{\hbar/(2T)} \). Er folgt aus der diskreten Zeit.
\end{itemize}

\section{Die Bohm-Potential Statistik}

Die zehnte Zeile zeigt die Statistik des Bohm-Potentials:

\begin{verbatim}
Q_mean: -2022.872            - Mittelwert des Bohm-Potentials Q
\end{verbatim}

Die Bedeutung:

\begin{itemize}
\item \textbf{\( Q_{\text{mean}} \):} Der Mittelwert von \( Q \). Er ist negativ – \( Q \) wirkt bindend.
\item \textbf{Negatives \( Q \):} \( Q \) ist die Quelle der Gravitation.
\end{itemize}

\section{Die physikalische Interpretation}

Die Werte der Simulation haben eine klare physikalische Bedeutung:

\begin{itemize}
\item \( \sum |I|^2 = \text{const} \) – Axiom 2 ist erfüllt. Die Information bleibt erhalten.
\item \( I_{\text{max}} > I_{\text{min}} \) – Struktur existiert. Materie ist da.
\item \( Q_{\text{mean}} < 0 \) – \( Q \) ist bindend. Gravitation ist da.
\item \( \vec{F} = -\nabla Q \) – Gravitation emergiert aus \( Q \).
\item \( J_{\text{mean}} \approx 0 \) – Das System ist im Gleichgewicht. Keine Expansion.
\end{itemize}

\section{Das Fazit der Simulation}

Die Simulation läuft stabil.

Axiom 2 ist erfüllt.

Strukturen existieren.

Das System ist im Gleichgewicht.

Alle Werte sind konsistent.

Die Simulation ist ein physischer Beweis, dass die IWT funktioniert.

\section{Zusammenfassung}

\begin{itemize}
\item Die Ausgabe der Simulation ist strukturiert und interpretierbar.
\item \( \sum |I|^2 \) ist konstant – Axiom 2 ist erfüllt.
\item \( I_{\text{max}} > I_{\text{min}} \) – Struktur existiert.
\item \( Q_{\text{mean}} < 0 \) – \( Q \) wirkt bindend.
\item \( \vec{F} = -\nabla Q \) – Gravitation emergiert aus \( Q \).
\item Die Simulation ist der Beweis, dass die IWT funktioniert.
\end{itemize}